\documentclass[10pt,a4paper]{amsart}
\usepackage[margin=19mm]{geometry}
\usepackage{amsmath,amssymb,mathtools}
\usepackage[scr=rsfs,cal=euler]{mathalfa}
\usepackage{xfrac}
\usepackage[unicode,hidelinks,bookmarks=false]{hyperref}
\newcommand{\ie}{i.e.\ }
\newcommand{\cf}{cf.\ }
\newcommand{\resp}{respectively\ }
\newcommand{\Subsection}{\subsection}
\providecommand{\nobreakdash}{\nobreak\hbox{-}}
\setlength{\emergencystretch}{2em}
\multlinegap0pt
\usepackage{bm}
\usepackage{bbm}
\usepackage{dsfont}
\usepackage{xspace}
\usepackage{cancel}
\usepackage{mathtools}
\usepackage{amsthm}
\makeatletter
\@ifundefined{thm}{\newtheorem{thm}{Theorem}}{}
\@ifundefined{lem}{\newtheorem{lem}[thm]{Lemma}}{}
\makeatother
\newcommand\ZZ{\mathbb{Z}}
\newcommand{\Aut}{\operatorname{Aut}}
\newcommand{\Inn}{\operatorname{Inn}}
\title[Twisted conjugacy in $BS(n, 1)$]{Twisted conjugacy in $BS(n, 1)$}
\author{Oorna Mitra, Mallika Roy, Enric Ventura}
\date{Source: arXiv:2508.04397v1 (CC BY 4.0)}
\begin{document}
\maketitle

\noindent\emph{Source and licence.} Twisted conjugacy in $BS(n, 1)$. Version arXiv:2508.04397v1 (CC BY 4.0), distributed under the Creative Commons Attribution 4.0 International licence, as recorded in the arXiv licence metadata for this exact version and shown on the abstract page https://arxiv.org/abs/2508.04397v1. Full rights notice: licenses/arxiv-2508.04397-CC-BY-4.0.txt.\par\smallskip

\noindent\textit{Editorial scope.} Counted unit: the Thm whose source label is \texttt{thm: main TCP} in the article (source0.tex lines 629--631, source label \texttt{thm: main TCP}), delivered together with the article's own complete original proof of it (source0.tex lines 633--696, one \texttt{proof} environment of 64 lines). No mathematics is written, repaired or re-derived: statement and proof are the article's own text line for line. Required context retained uncounted: eq: inverse (lines 398--400); eq: action (lines 478--480); lem: twisted autos (lines 620--627). Every definition, display and label of those lines is retained unchanged. Omitted from this package and not redistributed: 1--397, 401--477, 481--619, 628--628, 632--632, 697--1008. Nothing inside a retained excerpt is omitted or altered: every retained body is delivered exactly as the article's lines are written, so no omission receipt of any kind accompanies this package. Citations: the retained excerpts carry no citation command at all, so no bibliography entry of the article is transcribed and no reference is redistributed. Reference adaptation: none was needed and none was made; every reference inside the delivered unit resolves inside it, so no unresolved reference or citation is produced. Printed slips kept literal: no printed word, symbol, hypothesis or number of the counted statement or of its proof was altered. Layout and class: the article's own class is \texttt{amsart} with packages including amsfonts, amsmath, amssymb, amscd, amsthm, array, graphicx, inputenc, mathabx, stmaryrd, dsfont, babel, color, xy; the wrapper is the lane's pinned \texttt{amsart} editorial wrapper at 10pt on A4, which restates the theorem-like environments the retained text needs and adds the article's own macro definitions that the retained text needs (\texttt{\textbackslash Aut}, \texttt{\textbackslash Inn}, \texttt{\textbackslash ZZ}), with extra wrapper packages (bm, bbm, dsfont, xspace, cancel, mathtools). The delivered numbering is the unit's own. Assets: no retained span contains a figure environment, a graphics-inclusion command, a PDF-page inclusion, a table or a TikZ picture; no image file is copied into this package, the package declares no asset dependency and no image file is redistributed with it. Licence: version arXiv:2508.04397v1 of this article is distributed under the Creative Commons Attribution 4.0 International licence, as recorded in the arXiv licence metadata for this exact version and shown on the abstract page https://arxiv.org/abs/2508.04397v1. A full rights notice is delivered at licenses/arxiv-2508.04397-CC-BY-4.0.txt. Characters: every retained excerpt is pure ASCII, so the delivered engine, XeLaTeX with its default Latin Modern text font, reports no missing character.
 \begin{equation}\label{eq: inverse}
\big( a^{\alpha}t^c \big)^{-1} =t^{-c}a^{-\alpha}=a^{-n^{-c}\alpha}t^{-c}.
 \end{equation}

 \begin{align}\label{eq: action}
\big( a^{\nu}t^c \big) \varphi_{\alpha,\beta} = \big( t^{-r}a^{m}t^{r+c}\big) \varphi_{\alpha, \beta} &= \big( a^{\beta}t\big)^{-r} a^{m\alpha} \big(a^{\beta}t\big)^{r+c}= \notag \\ &= a^{\frac{n^{-r}-1}{n-1}\beta}t^{-r} a^{m\alpha} a^{\frac{n^{r+c}-1}{n-1}\beta}t^{r+c}= \notag \\ &= a^{\frac{n^{-r}-1}{n-1}\beta} a^{m n^{-r}\alpha} a^{n^{-r}\frac{n^{r+c}-1}{n-1}\beta}t^{c}= \\ &= a^{\frac{n^{-r}-1}{n-1}\beta+ m n^{-r}\alpha + \frac{n^{c}-n^{-r}}{n-1}\beta}t^{c}= \notag \\ &= a^{\nu \alpha + \frac{n^{c}-1}{n-1}\beta}t^{c}. \notag
 \end{align}

\begin{lem} \label{lem: twisted autos}
Let $G$ be a group, $u,v,w \in G$, $\varphi, \psi \in \Aut(G)$ and $\gamma_w \in \Inn(G)$. Then the following are equivalent:
 \begin{itemize}
\item[(i)] $u \sim_\phi v$,
\item[(ii)] $u \psi \sim_{\psi^{-1}\varphi\psi} v \psi$,
\item[(iii)] $wu \sim_{\varphi\gamma_{w^{-1}}} wv$, \item[(iv)] $uw \sim_{\gamma_{w^{-1}}\varphi} vw$. \hfill \qed
 \end{itemize}
\end{lem}

\begin{thm}\label{thm: main TCP}
The Twisted Conjugacy Problem is solvable in $BS(n,1)$.
\end{thm}

\begin{proof}
Suppose we are given a fixed automorphism $\varphi_{\alpha, \beta}\in \Aut(BS(n,1))$, 
%  $$
% \begin{array}{rcl} \varphi_{\alpha, \beta}\colon BS(n,1) & \to & BS(n,1) \\ a & \mapsto & a^{\alpha} \\ t & \mapsto & a^{\beta} t,
%  \end{array}
%  $$
where $\alpha=k/n^p \in \ZZ[1/n]^*, \beta=\ell/n^q\in \ZZ[1/n]$ with $k,\ell, p,q\in \ZZ$, $p,q\geq 0$. Suppose we are also given two elements $u=a^{\nu_1}t^{c_1},\, v=a^{\nu_2}t^{c_2}\in BS(n,1)$, where $\nu_i =m_i/n^{r_i}\in \ZZ[1/n]$ with $m_i, r_i, c_i\in \ZZ$, $r_i\geq 0$, $i=1,2$. We have to decide whether $a^{\nu_1}t^{c_1}$ and $a^{\nu_2}t^{c_2}$ are $\varphi_{\alpha, \beta}$-twisted conjugated to each other, i.e., whether $u\sim_{\varphi_{\alpha, \beta}} v$. Observe that, again, $c_1=c_2$ (call it just $c$) is an obvious necessary condition. 

Applying Lemma~\ref{lem: twisted autos} (i)$\Leftrightarrow$(iv) with $w=t^{-c}$, we can assume $c_1=c_2=0$, i.e., $a^{\nu_i}t^{c_i}=t^{-r_i}a^{m_i}t^{r_i}$, $i=1,2$. Then, applying Lemma~\ref{lem: twisted autos} (i)$\Leftrightarrow$(ii) with $\psi=\gamma_{t^{-r_1}}$, we can assume $r_1=0$. In other words, we are reduced to inputs of the form $u=a^{m_1},\, v=a^{\nu_2}=t^{-r_2}a^{m_2}t^{r_2}\in BS(n,1)$.

Writing the possible $\varphi_{\alpha, \beta}$-twisted conjugator as $g=a^{\nu}t^d$, $\nu\in \ZZ[1/n]$, $d\in \ZZ$, and using equations~\eqref{eq: action} and~\eqref{eq: inverse}, we have 
 $$
g\varphi_{\alpha,\beta} =a^{\nu \alpha + \frac{n^{d}-1}{n-1}\beta}t^{d},
 $$
 $$
\big( g\varphi_{\alpha,\beta} \big)^{-1} =a^{-n^{-d}(\nu \alpha + \frac{n^{d}-1}{n-1}\beta)}t^{-d},
 $$
and 
 $$
\big( g\varphi_{\alpha,\beta} \big)^{-1} ug= a^{-n^{-d}(\nu \alpha + \frac{n^{d}-1}{n-1}\beta)} t^{-d} a^{m_1} a^{\nu}t^d= a^{-n^{-d}\nu \alpha + \frac{n^{-d}-1}{n-1}\beta} a^{m_1 n^{-d}} a^{\nu n^{-d}}.
 $$
So, $a^{m_1}\sim_{\varphi_{\alpha, \beta}} a^{\nu_2}$ if and only if 
 $$
\frac{m_2}{n^{r_2}} n^d =-\nu \alpha +\frac{1-n^{d}}{n-1}\beta +m_1 +\nu, 
 $$
 $$
\Leftrightarrow (n-1)m_2 n^d +(n-1)n^{r_2}\nu (\alpha -1) +(n^{d}-1)n^{r_2}\beta =(n-1)n^{r_2}m_1, 
 $$
for some $\nu\in \ZZ[1/n]$ and $d\in \ZZ$. Writing $\nu =y/n^x$, $x,y\in \ZZ$, $x\geq 0$, and applying the change of variable $\ZZ \ni z=d+x$, this is equivalent to the integral equation  
 $$
(n-1)m_2 n^d +(n-1)n^{r_2}\frac{y}{n^x} \frac{k-n^p}{n^p} +(n^{d}-1)n^{r_2}\frac{\ell}{n^q} =(n-1)n^{r_2}m_1, 
 $$
 $$
\Rightarrow (n-1)m_2 n^{d+x+p+q} +(n-1)y(k-n^p)n^{r_2+q} +(n^{d}-1)\ell n^{r_2+p+x} =(n-1)m_1 n^{r_2+p+q+x}, 
 $$
 $$
\Rightarrow \Big( \ell n^{r_2+p}+(n-1)m_1 n^{r_2+p+q}\Big) n^x +\Big( (n-1)(n^p-k)n^{r_2+q} \Big) y =\Big( (n-1)m_2 n^{p+q} +\ell n^{r_2+p}\Big) n^{z}.
 $$
From the data $k, p, \ell, q, m_1, m_2, r_2, n\in \ZZ$, $p,q,r_2,n\geq 0$, compute the three integers 
 $$
A=\ell n^{r_2+p}+(n-1)m_1 n^{r_2+p+q},\qquad B=(n-1)(n^p-k)n^{r_2+q}, \qquad C=(n-1)m_2 n^{p+q}+\ell n^{r_2+p}.
 $$
Summarizing, $A,B,C\in \ZZ$, and $u$ and $v$ are $\varphi_{\alpha, \beta}$-twisted conjugated to each other if and only if the equation 
 \begin{equation}\label{eq: equation}
An^x +By=Cn^z
 \end{equation}
has an integral solution, for the unknowns $x,y,z\in \ZZ$, with $x\geq 0$. Write $C=n^s C'$ with $s\geq 0$ and $C'$ not being multiple of $n$, and note that all possible solutions to equation~\eqref{eq: equation} have $z\geq -s$. So, with the change of variable $\ZZ \ni z'=z+s$, \eqref{eq: equation} is equivalent to 
 \begin{equation}\label{eq: equation2}
An^x +By=C'n^{z'},
 \end{equation}
for the unknowns $x,y,z'\in \ZZ$, with $x,z'\geq 0$. Let us distinguish three cases. 

\textbf{Case 1: $B=\pm 1$}. Obviously, equation~\eqref{eq: equation2} always has solutions.  

\textbf{Case 2: $B=0$}. Our equation becomes $An^x=C' n^{z'}$, which has the desired solution if and only if either $A=C'=0$, or both $A,C'\neq 0$ and $A/C$ is a rational number being a (positive or negative) power of $n$. This can be easily checked by looking at the prime factorizations of $A$, $C'$, and $n$.

\textbf{Case 3: $B\neq 0$}. Now, equation~\eqref{eq: equation2} is equivalent to 
 $$
An^x \equiv C'n^{z'} \pmod{B},
 $$
for the unknowns $x,z'\in \ZZ$, with $x,z'\geq 0$. Compute the finite set $\mathcal{N}=\{1,n,n^2, \ldots \}\subseteq \ZZ/B\ZZ$ (by computing the successive powers of $n$ until obtaining the first repetition modulo $B$). Clearly, our equation admits a solution of and only if the subsets $A\mathcal{N}$ and $C'\mathcal{N}$ from $\ZZ/B\ZZ$ intersect non-trivially. 

This completes the proof.
\end{proof}

\end{document}
